\documentclass[11pt,a4paper]{article}
\usepackage[margin=1in]{geometry}
\usepackage{amsmath,amssymb,amsthm}
\usepackage{algorithm}
\usepackage{algpseudocode}
\usepackage{booktabs}
\usepackage{hyperref}

\theoremstyle{definition}
\newtheorem{definition}{\textbf{Definition}}[section]
\newtheorem{theorem}{\textbf{Theorem}}[section]
\newtheorem{lemma}[theorem]{\textbf{Lemma}}
\newtheorem{remark}{\textbf{Remark}}[section]

\title{\textbf{Generative Adversarial Networks}}
\author{\textbf{Team Scaibu} \\ \small Email: \texttt{jaydeep.9664920749@gmail.com} \\ \small \textbf{Architecture and Core Engine Team}}
\date{\textbf{October 2026}}

\begin{document}
\maketitle

\section{Definitions and Cognitive Mental Models}

\subsection{Formal Mathematical Definition}
\textbf{Definition (Generative Adversarial Game Operator):}
Let $\mathcal{X} = \mathbb{R}^{d_x}$ denote the ambient data space and $\mathcal{Z} = \mathbb{R}^{d_z}$ denote the latent noise prior space with density $p_z(\mathbf{z})$. Let $G_\theta: \mathcal{Z} \to \mathcal{X}$ be a differentiable parameterized generator mapping noise $\mathbf{z}$ to synthetic samples, inducing an implicit probability density $p_g$. Let $D_\phi: \mathcal{X} \to (0, 1)$ denote a differentiable discriminator predicting the probability that a sample originates from the true data distribution $p_{\text{data}}$ rather than $p_g$.

The zero-sum game is formalized as the two-player minimax optimization:
{\boldmath
\begin{equation}
\min_G \max_D V(D, G) = \mathbb{E}_{\mathbf{x} \sim p_{\text{data}}}[\log D(\mathbf{x})] + \mathbb{E}_{\mathbf{z} \sim p_z}[\log(1 - D(G(\mathbf{z})))]
\end{equation}
}
In standard deep network training, to prevent gradient saturation early in optimization when $D$ easily rejects poor generator samples, the generator minimizes the \textbf{non-saturating objective}:
{\boldmath
\begin{equation}
\mathcal{L}_G^{\text{NS}} = -\mathbb{E}_{\mathbf{z} \sim p_z}[\log D(G(\mathbf{z}))]
\end{equation}
}

\subsection{Expanded Non-Technical Definition (Intuition for Anyone)}
\textbf{Intuitive Conceptual Definition:}
Most machine learning systems are trained like students preparing for a multiple-choice test: a teacher marks their answers right or wrong according to a fixed answer key. Generative Adversarial Networks (GANs) discard the answer key and instead set up a high-stakes, competitive simulation between two opposing players:
\begin{enumerate}
    \item \textbf{The Forger (The Generator)}: A student who has never seen real masterpieces, but tries to forge one out of random noise.
    \item \textbf{The Detective (The Discriminator)}: An investigator who examines paintings and tries to guess whether each one is a genuine museum masterpiece or a counterfeit made by the forger.
    \item \textbf{The Feedback Loop}: When the detective spots a flaw (e.g., \emph{``The brushstrokes are too blurry!''}), the forger immediately learns from that failure and adjusts its technique to produce sharper strokes.
    \item \textbf{The Nash Equilibrium}: Over thousands of rounds of competition, the forger becomes so extraordinarily skilled that its creations are physically indistinguishable from genuine masterpieces. At that point, the detective can only guess with 50\% accuracy (a coin flip).
\end{enumerate}

\subsection{Expanded Mental Model for Visualization (Understandable by a Child)}
\textbf{The Toy Forger and the Detective Bear}:
Imagine a playful competition in a toy workshop:
\begin{itemize}
    \item \textbf{The Genuine Golden Teddy Bear ($\mathbf{x} \sim p_{\text{data}}$)}:
    In the town square sits a real, handcrafted golden teddy bear with soft ears, button eyes, and velvet paws.
    
    \item \textbf{The Clay Modeler (Generator $G$)}:
    A young apprentice modeler is given a lump of gray clay and random dice ($\mathbf{z}$). Without ever directly touching the real bear, the apprentice tries to mold a look-alike bear.
    
    \item \textbf{The Detective Bear (Discriminator $D$)}:
    A grumpy detective bear examines each teddy bear. If it sees a real bear, it shouts: \emph{``Real! (100\%)''}. If it sees a lumpy gray clay bear, it shouts: \emph{``Fake! (0\%)''}.
    
    \item \textbf{The Helpful Secret Clue (Gradient $\nabla_D$)}:
    Every time the detective rejects a fake bear, he explains why: \emph{``The ears are too pointy!''} The clay modeler listens carefully and smooths the ears down.
    
    \item \textbf{The Perfect Match (Nash Equilibrium)}:
    After practicing hundreds of times, the apprentice molds a bear so soft, golden, and perfect that the detective bear rubs his eyes in confusion: \emph{``I can't tell the difference! It's 50/50!''}
\end{itemize}

\section{Core Mathematical Equations: Formal Definitions, Meaning, and Importance}

\subsection{\textbf{Equation 1: Minimax Cross-Entropy Objective}}
{\boldmath
\begin{equation}
V(D, G) = \mathbb{E}_{\mathbf{x} \sim p_{\text{data}}}[\log D(\mathbf{x})] + \mathbb{E}_{\mathbf{x} \sim p_g}[\log(1 - D(\mathbf{x}))]
\end{equation}
}
\begin{itemize}
    \item \textbf{Formal Mathematical Definition}:
    The expected payoff functional of a two-player zero-sum game, corresponding to the negative cross-entropy loss of binary classification between distributions $p_{\text{data}}$ and $p_g$.
    \item \textbf{What It Means}:
    Measures how effectively the discriminator separates real and generated samples.
    \item \textbf{Why It Is Important}:
    Provides the foundational mathematical framework proving that minimax optimization minimizes divergence between empirical and model distributions.
\end{itemize}

\subsection{\textbf{Equation 2: Optimal Bayes Discriminator Solution}}
{\boldmath
\begin{equation}
D^*(\mathbf{x}) = \frac{p_{\text{data}}(\mathbf{x})}{p_{\text{data}}(\mathbf{x}) + p_g(\mathbf{x})}
\end{equation}
}
\begin{itemize}
    \item \textbf{Formal Mathematical Definition}:
    The functional maximizer of $V(D, G)$ for any fixed generator $G$, obtained by variational differentiation $\frac{\delta V}{\delta D} = 0$.
    \item \textbf{What It Means}:
    At optimality, the discriminator outputs the exact posterior probability that sample $\mathbf{x}$ came from the real data distribution under equal priors.
    \item \textbf{Why It Is Important}:
    Allows substituting $D^*$ back into the minimax game to derive the theoretical connection to classical probability divergences.
\end{itemize}

\subsection{\textbf{Equation 3: Jensen-Shannon Divergence Equivalence}}
{\boldmath
\begin{equation}
V(D^*, G) = -\log(4) + 2 \cdot D_{\text{JS}}(p_{\text{data}} \parallel p_g)
\end{equation}
}
\begin{itemize}
    \item \textbf{Formal Mathematical Definition}:
    The value of the minimax game evaluated at the optimal discriminator, expressing the objective as an affine transformation of the symmetrized Jensen-Shannon divergence $D_{\text{JS}}(p \parallel q) = \frac{1}{2} D_{\text{KL}}(p \parallel \frac{p+q}{2}) + \frac{1}{2} D_{\text{KL}}(q \parallel \frac{p+q}{2})$.
    \item \textbf{What It Means}:
    Minimizing the generator loss against an optimal discriminator is strictly equivalent to minimizing the Jensen-Shannon divergence between the real and generated data distributions.
    \item \textbf{Why It Is Important}:
    Guarantees that the global minimum of the game is unique and attained if and only if $p_g = p_{\text{data}}$, where $V(D^*, G) = -\log(4)$ and $D^*(\mathbf{x}) = 1/2$.
\end{itemize}

\subsection{\textbf{Equation 4: Non-Saturating Generator Gradient}}
{\boldmath
\begin{equation}
\mathcal{L}_G^{\text{NS}} = -\mathbb{E}_{\mathbf{z}}[\log D(G(\mathbf{z}))] \quad \Longrightarrow \quad \nabla_\theta \mathcal{L}_G^{\text{NS}} = -\mathbb{E}_{\mathbf{z}}\left[ \frac{\nabla_\theta D(G(\mathbf{z}))}{D(G(\mathbf{z}))} \right]
\end{equation}
}
\begin{itemize}
    \item \textbf{Formal Mathematical Definition}:
    The surrogate maximization objective for the generator that replaces $\log(1 - D(G(\mathbf{z})))$ with $-\log D(G(\mathbf{z}))$.
    \item \textbf{What It Means}:
    When the generator is weak and $D(G(\mathbf{z})) \approx 0$, the original objective has gradient $\approx 0$ (saturation). The non-saturating objective provides large, non-vanishing gradient signals early in training.
    \item \textbf{Why It Is Important}:
    Rescues GAN optimization from catastrophic early gradient vanishing, making deep adversarial training feasible in practice.
\end{itemize}

\section{Rigorous Mathematical Analysis Checklist}

\subsection{[ ] Notation: Symbol and Operator Specification}
\begin{itemize}
    \item $\mathbf{x} \in \mathbb{R}^{d_x}$: Real data vector sampled from $p_{\text{data}}$.
    \item $\mathbf{z} \in \mathbb{R}^{d_z}$: Latent noise vector sampled from prior $p_z$.
    \item $G(\mathbf{z}) \in \mathbb{R}^{d_x}$: Generator output sample.
    \item $D(\mathbf{x}) \in (0, 1)$: Discriminator scalar realness probability.
    \item $D_{\text{JS}}$: Jensen-Shannon divergence.
    \item $\alpha \in [0, 0.5)$: One-sided label smoothing factor.
\end{itemize}

\subsection{[ ] Definitions: Epistemic Nature of Variables}
\begin{table}[h]
\centering
\caption{\textbf{Epistemic Classification of GAN Quantities}}
\begin{tabular}{lll}
\toprule
\textbf{Variable} & \textbf{Epistemic Classification} & \textbf{Mathematical Nature} \\
\midrule
$d_x, d_z$ & \textbf{Defined} (Architectural hyperparameter) & Natural numbers $\mathbb{N}_{\ge 1}$ \\
$\alpha$ & \textbf{Defined} (Label smoothing constant) & Non-negative real scalar in $[0, 0.5)$ \\
$G_\theta, D_\phi$ & \textbf{Learned} (Adversarial neural networks) & Parameterized non-linear tensor mappings \\
$D(\mathbf{x}), D(G(\mathbf{z}))$ & \textbf{Calculated} (Discriminator predictions) & Real probabilities in $(0, 1)$ \\
$\mathcal{L}_D, \mathcal{L}_G$ & \textbf{Calculated} (Cross-entropy losses) & Real optimization objective scalars \\
\bottomrule
\end{tabular}
\end{table}

\subsection{[ ] Operator Computation Graph}
{\boldmath
\begin{equation}
\mathbf{z} \xrightarrow{G} \mathbf{x}_{\text{fake}} \xrightarrow{D} D(\mathbf{x}_{\text{fake}}) \longrightarrow \mathcal{L}_G, \qquad (\mathbf{x}_{\text{real}}, \mathbf{x}_{\text{fake}}) \xrightarrow{D} (D_{\text{real}}, D_{\text{fake}}) \longrightarrow \mathcal{L}_D
\end{equation}
}

\subsection{[ ] Dimensional Units and Tensor Ranks}
\begin{itemize}
    \item Noise $\mathbf{z}$: Rank-1 tensor of shape $(d_z)$.
    \item Samples $\mathbf{x}, G(\mathbf{z})$: Rank-1 tensor of shape $(d_x)$.
    \item Discriminator outputs $D(\mathbf{x})$: Dimensionless probabilities in $[0, 1]$.
    \item Losses $\mathcal{L}_D, \mathcal{L}_G$: Dimensionless scalar loss values.
\end{itemize}

\subsection{[ ] Limiting Regimes and Asymptotic Behaviors}
\begin{itemize}
    \item \textbf{Optimal Discriminator Limit ($D \to D^*$)}: Generator gradient aligns precisely with the gradient of Jensen-Shannon divergence.
    \item \textbf{Disjoint Supports Limit ($\operatorname{supp}(p_{\text{data}}) \cap \operatorname{supp}(p_g) = \emptyset$)}: $D_{\text{JS}} \equiv \log(2)$ is constant; discriminator easily achieves 100\% accuracy and gradients vanish everywhere.
\end{itemize}

\subsection{[ ] Operational Constraints and Failure Modes}
\begin{itemize}
    \item \textbf{Mode Collapse}: Generator maps all noise inputs $\mathbf{z}$ to a small subset of realistic samples, ignoring the full diversity of the true distribution.
    \item \textbf{Vanishing Gradients}: If discriminator trains too quickly to near-zero error, $\nabla_{\mathbf{x}} \log(1 - D(G(\mathbf{z}))) \to \mathbf{0}$, freezing the generator.
\end{itemize}

\subsection{[ ] Mathematical Invariants and Conserved Quantities}
\begin{itemize}
    \item \textbf{Nash Equilibrium Invariant}: At convergence, $p_g = p_{\text{data}}$, $D^*(\mathbf{x}) = 0.5$ identically for all $\mathbf{x} \in \operatorname{supp}(p_{\text{data}})$.
    \item \textbf{Symmetry of JSD}: $D_{\text{JS}}(p \parallel q) = D_{\text{JS}}(q \parallel p) \in [0, \log 2]$.
\end{itemize}

\subsection{[ ] Cross-Domain Mathematical Isomorphisms}
\begin{itemize}
    \item \textbf{Zero-Sum Game Theory}: Direct instantiation of von Neumann's minimax theorem in infinite-dimensional probability function space.
    \item \textbf{F-Divergence Minimization}: GAN training is a specific variational representation of general $f$-divergence estimation.
\end{itemize}

\section{Theoretical Theorems and Formal Proofs}

\begin{theorem}[\textbf{Global Optimality of the Minimax Game}]
For any fixed generator $G$, the optimal discriminator is $D^*_G(\mathbf{x}) = \frac{p_{\text{data}}(\mathbf{x})}{p_{\text{data}}(\mathbf{x}) + p_g(\mathbf{x})}$. Furthermore, the infimum $\min_G V(D^*_G, G)$ is uniquely attained if and only if $p_g = p_{\text{data}}$, where $V(D^*_G, G) = -\log 4$.
\end{theorem}

\begin{proof}
For any given generator $G$, the discriminator objective is:
{\boldmath
\begin{align}
V(D, G) &= \int_{\mathcal{X}} p_{\text{data}}(\mathbf{x}) \log D(\mathbf{x}) \, d\mathbf{x} + \int_{\mathcal{Z}} p_z(\mathbf{z}) \log(1 - D(G(\mathbf{z}))) \, d\mathbf{z} \\
&= \int_{\mathcal{X}} \left[ p_{\text{data}}(\mathbf{x}) \log D(\mathbf{x}) + p_g(\mathbf{x}) \log(1 - D(\mathbf{x})) \right] d\mathbf{x}
\end{align}
}
For any fixed $\mathbf{x}$, consider maximizing the integrand $f(y) = a \log y + b \log(1 - y)$ over $y \in (0, 1)$, where $a = p_{\text{data}}(\mathbf{x})$ and $b = p_g(\mathbf{x})$. Differentiating with respect to $y$:
{\boldmath
\begin{equation}
f'(y) = \frac{a}{y} - \frac{b}{1 - y} = 0 \quad \Longrightarrow \quad a(1 - y) = b y \quad \Longrightarrow \quad y = \frac{a}{a + b}
\end{equation}
}
Since $f''(y) = -\frac{a}{y^2} - \frac{b}{(1 - y)^2} < 0$, this critical point is a global maximum:
{\boldmath
\begin{equation}
D^*_G(\mathbf{x}) = \frac{p_{\text{data}}(\mathbf{x})}{p_{\text{data}}(\mathbf{x}) + p_g(\mathbf{x})}
\end{equation}
}
Substituting $D^*_G(\mathbf{x})$ back into $V(D, G)$:
{\boldmath
\begin{align}
V(D^*_G, G) &= \int_{\mathcal{X}} \left[ p_{\text{data}}(\mathbf{x}) \log \frac{p_{\text{data}}(\mathbf{x})}{p_{\text{data}}(\mathbf{x}) + p_g(\mathbf{x})} + p_g(\mathbf{x}) \log \frac{p_g(\mathbf{x})}{p_{\text{data}}(\mathbf{x}) + p_g(\mathbf{x})} \right] d\mathbf{x} \\
&= -\log 4 + \int p_{\text{data}}(\mathbf{x}) \log \frac{p_{\text{data}}(\mathbf{x})}{\frac{p_{\text{data}}(\mathbf{x}) + p_g(\mathbf{x})}{2}} d\mathbf{x} + \int p_g(\mathbf{x}) \log \frac{p_g(\mathbf{x})}{\frac{p_{\text{data}}(\mathbf{x}) + p_g(\mathbf{x})}{2}} d\mathbf{x} \\
&= -\log 4 + 2 \cdot D_{\text{JS}}(p_{\text{data}} \parallel p_g)
\end{align}
}
Since Jensen-Shannon divergence is strictly non-negative and vanishes if and only if $p_g = p_{\text{data}}$, the minimum value is $-\log 4$, achieved uniquely when $p_g = p_{\text{data}}$.
\end{proof}

\section{Foundational Architectural and Epistemic Meta-Questions}

\subsection{Architectural Question 1: Why Does JSD Cause Mode Collapse in Practice?}
\textbf{Question}: If the theoretical minimum is attained when $p_g = p_{\text{data}}$, why do practical GANs frequently collapse to a few repetitive modes?

\textbf{Meta-Question}: Is the optimization topology of a minimax objective fundamentally different from single-objective gradient descent?

\textbf{Answer}: In practice, real data manifolds lie on thin, low-dimensional subsets of $\mathbb{R}^{d_x}$. When $p_g$ and $p_{\text{data}}$ do not share support, $D_{\text{JS}}(p_{\text{data}} \parallel p_g) = \log 2$ everywhere, producing zero gradient across most of the ambient space. Furthermore, standard gradient descent-ascent on $\min_G \max_D V$ is not guaranteed to converge to a Nash equilibrium; it can exhibit rotational limit cycles. The generator finds it easiest to fool the discriminator by concentrating all its probability mass on a single high-scoring mode (mode collapse) rather than covering the complex full data distribution.

\subsection{Architectural Question 2: Why Is One-Sided Label Smoothing Stabilizing?}
\textbf{Question}: Why should label smoothing be applied only to real labels ($1.0 \to 1.0 - \alpha$) and never to fake labels ($0.0 \to \alpha$)?

\textbf{Meta-Question}: How does asymmetric regularization prevent deceptive training gradients?

\textbf{Answer}: If fake labels are smoothed to $\alpha > 0$, the optimal discriminator becomes $D^*(\mathbf{x}) = \frac{(1-\alpha)p_{\text{data}} + \alpha p_g}{p_{\text{data}} + p_g}$. Notice that in regions where $p_{\text{data}}(\mathbf{x}) = 0$ but $p_g(\mathbf{x}) > 0$, the optimal output is $D^*(\mathbf{x}) = \alpha > 0$. The discriminator is therefore explicitly incentivized to reinforce generator samples in regions with zero real data probability. Smoothing only real targets preserves the zero-penalty condition on fake samples while preventing discriminator weights from blowing up toward infinity on real samples.

\section{Algorithmic Specification}

\begin{algorithm}[H]
\caption{Generative Adversarial Network Forward Loss Computation}
\label{alg:gan}
\begin{algorithmic}[1]
\Require Real probabilities $\mathbf{p}_{\text{real}} \in (0, 1)^{N_r}$, Fake probabilities $\mathbf{p}_{\text{fake}} \in (0, 1)^{N_f}$, Smoothing $\alpha \ge 0$
\Ensure Discriminator loss $\mathcal{L}_D$, Generator loss $\mathcal{L}_G$, Real accuracy $A_r$, Fake accuracy $A_f$
\State $\epsilon \gets 10^{-12}, \quad y_{\text{real}} \gets 1.0 - \alpha$
\State $L_{\text{real\_sum}} \gets 0.0, \quad C_r \gets 0$
\For{$i \gets 1 \textbf{ to } N_r$}
    \State $p \gets \min(1.0 - \epsilon, \max(\epsilon, p_{\text{real}}[i]))$
    \State $L_{\text{real\_sum}} \gets L_{\text{real\_sum}} - (y_{\text{real}} \cdot \ln(p) + (1.0 - y_{\text{real}}) \cdot \ln(1.0 - p))$
    \If{$p_{\text{real}}[i] \ge 0.5$}
        \State $C_r \gets C_r + 1$
    \EndIf
\EndFor
\State $L_{\text{fake\_sum}} \gets 0.0, \quad L_{G\text{\_sum}} \gets 0.0, \quad C_f \gets 0$
\For{$j \gets 1 \textbf{ to } N_f$}
    \State $p \gets \min(1.0 - \epsilon, \max(\epsilon, p_{\text{fake}}[j]))$
    \State $L_{\text{fake\_sum}} \gets L_{\text{fake\_sum}} - \ln(1.0 - p)$
    \State $L_{G\text{\_sum}} \gets L_{G\text{\_sum}} - \ln(p)$ \Comment{Non-saturating generator loss}
    \If{$p_{\text{fake}}[j] < 0.5$}
        \State $C_f \gets C_f + 1$
    \EndIf
\EndFor
\State $\mathcal{L}_D \gets 0.5 \cdot (L_{\text{real\_sum}} / N_r + L_{\text{fake\_sum}} / N_f)$
\State $\mathcal{L}_G \gets L_{G\text{\_sum}} / N_f$
\State $A_r \gets C_r / N_r, \quad A_f \gets C_f / N_f$
\State \Return $\{\text{discriminator\_loss}: \mathcal{L}_D, \text{generator\_loss}: \mathcal{L}_G, \text{d\_real\_acc}: A_r, \text{d\_fake\_acc}: A_f\}$
\end{algorithmic}
\end{algorithm}

\section{Complexity Analysis}
\begin{itemize}
    \item \textbf{Time Complexity}:
    \begin{itemize}
        \item Real loss evaluation: $N_r$ clamping, log, and accumulation operations ($\Theta(N_r)$).
        \item Fake and generator loss evaluation: $N_f$ clamping, log, and accumulation operations ($\Theta(N_f)$).
        \item Total Time: $\Theta(N_r + N_f)$ FLOPs.
    \end{itemize}
    \item \textbf{Space Complexity}:
    \begin{itemize}
        \item Auxiliary memory: $\Theta(1)$ scalar accumulators.
        \item Total Space: $\Theta(1)$.
    \end{itemize}
\end{itemize}

\section{Repository Reference}
All mathematical formulations, algorithmic implementations, contracts, and test suites are maintained at:
\begin{center}
\url{https://github.com/Chief-Strategist-J/llm-obs-infra}
\end{center}

\end{document}
